% Convexity of the local ball contains the full Newton-step segment.
\begin{tikzpicture}[x=1cm,y=1cm,font=\small,>=Latex]
\draw[black] (0,0) circle (1.35);
\node[above] at (0,1.45) {$B(\mathbf x^*,r)$};
\fill (0,0) circle (2pt) node[below=3mm] {$\mathbf x^*$};
\draw[very thick,blue!70!black] (.65,.3)--(-.65,.15);
\fill (.65,.3) circle (2pt) node[above right] {$\mathbf x^k$};
\fill (-.65,.15) circle (2pt) node[above=2mm] {$\mathbf x^{k+1}$};
\fill[orange!80!black] (.039,.2295) circle (2.5pt);
\begin{scope}[shift={(3.2,0)}]
\draw[very thick,blue!70!black] (0,0)--(7,0);
\fill (0,0) circle (2pt) node[below=2mm] {$\mathbf x^k$};
\fill (7,0) circle (2pt) node[below=2mm] {$\mathbf x^{k+1}$};
\fill[orange!80!black] (3.29,0) circle (2.5pt) node[below=5mm] {$\mathbf x^k+t\mathbf p^k$};
\draw[<->] (0,.55)--(3.29,.55) node[midway,above] {$t\|\mathbf p^k\|$};
\draw[<->,densely dashed] (0,1.3)--(7,1.3) node[midway,above] {$\|\mathbf p^k\|$};
\node[below] at (3.5,-1.05) {$0\le t\le1$; both endpoints and the segment lie in $B$.};
\end{scope}
\end{tikzpicture}
